\chapter{Independent line workers}\label{module:web-lineworkers}
\Q{Evidence}{Current implementation boundary}
\path{Software/DemoCode/web/lineworkers.py}. Exact file hashes are in the accompanying source manifest. Names, declarations and source locations below come from the captured source, not a runtime import.
\Q{Purpose}{Input-to-output reading flow}
Configure inherited readers/classifiers; decide whether pooling is eligible; initialise per-process thread limits; map independent line jobs and preserve order, with serial fallback. The server remains responsible for page-level orchestration.


\Q{Evidence}{Selected current module settings}
These are declarations or expressions, not recovered environment values. Source comments do not establish a measured optimum. Larger glyph coordinate tables are designed data in the bundled source.
\begin{lstlisting}[style=plain]
LINE_WORKERS = int(os.environ.get('LINE_WORKERS', '0'))
MIN_LINES_FOR_POOL = 4
_CFG = {}
\end{lstlisting}
\section*{Stage: Independent line workers}
\section{Method background: Independent line workers}\label{method:workers}
\Q{Purpose}{The role in this system}
Run independent text or writer passes on multiple processes while restoring page order.


\Q{Justification}{The reason for this construction}
Explicit fork workers inherit configured models and cap numerical threads to reduce oversubscription. The serial path remains available; startup and IPC still cost time.


\Q{Alternatives}{What changes if another method is used}
Serial processing avoids pool overhead for short pages. Threads share memory but interact with Python and BLAS scheduling. Spawn workers need explicit model initialisation rather than this fork-inherited configuration.


\Q{Library}{Implementation speed and accuracy}
The code already uses Python multiprocessing.Pool and NumPy/BLAS. Fork is copy-on-write, not a guarantee of no allocation; passing crops/results still needs communication. Forking a threaded server with native numerical state is a deployment consideration. No new speedup or identical-output test was run. \cite{pythonmp,numpymatmul}.


\section{Configure}\label{fn:web-lineworkers:Configure}
\begin{lstlisting}[style=sig]
def Configure(reader, namer, strokeNormalise)
\end{lstlisting}
\Q{Evidence}{Source and method link}
\path{Software/DemoCode/web/lineworkers.py} line 30. Implementation evidence: \cite{code2026}. Related method, alternatives and library discussion: Section~\ref{method:workers}.


\Q{Purpose}{What this contributes}
Run independent text or writer passes on multiple processes while restoring page order.


\Q{Method}{Current implementation}
Run independent text or writer passes on multiple processes while restoring page order. The excerpt below is from the current definition; the complete function is in the source snapshot. Its constants and branches, rather than an older signature, govern this path.


\begin{lstlisting}[style=plain]
def Configure(reader, namer, strokeNormalise):
    _CFG["reader"] = reader
    _CFG["namer"] = namer
    _CFG["norm"] = bool(strokeNormalise)
\end{lstlisting}
\Q{Justification}{Interpret the method rather than infer a benchmark}
Explicit fork workers inherit configured models and cap numerical threads to reduce oversubscription. The serial path remains available; startup and IPC still cost time. \cite{pythonmp,numpymatmul}.


\section{CanPool}\label{fn:web-lineworkers:CanPool}
\begin{lstlisting}[style=sig]
def CanPool(n)
\end{lstlisting}
\Q{Evidence}{Source and method link}
\path{Software/DemoCode/web/lineworkers.py} line 37. Implementation evidence: \cite{code2026}. Related method, alternatives and library discussion: Section~\ref{method:workers}.


\Q{Purpose}{What this contributes}
Allows a pool only when LINE\_WORKERS is not one, at least MIN\_LINES\_FOR\_POOL jobs are present, os.fork exists and configuration is nonempty. Eligibility uses the number of jobs passed to Map: after orientation probes this can be smaller than page line count.


\Q{Method}{Current implementation}
Run independent text or writer passes on multiple processes while restoring page order. The excerpt below is from the current definition; the complete function is in the source snapshot. Its constants and branches, rather than an older signature, govern this path.


\begin{lstlisting}[style=plain]
def CanPool(n):
    return (LINE_WORKERS != 1 and n >= MIN_LINES_FOR_POOL
            and hasattr(os, "fork") and bool(_CFG))
\end{lstlisting}
\Q{Justification}{Interpret the method rather than infer a benchmark}
Explicit fork workers inherit configured models and cap numerical threads to reduce oversubscription. The serial path remains available; startup and IPC still cost time. \cite{pythonmp,numpymatmul}.


\section{\_Workers}\label{fn:web-lineworkers:-Workers}
\begin{lstlisting}[style=sig]
def _Workers(n)
\end{lstlisting}
\Q{Evidence}{Source and method link}
\path{Software/DemoCode/web/lineworkers.py} line 43. Implementation evidence: \cite{code2026}. Related method, alternatives and library discussion: Section~\ref{method:workers}.


\Q{Purpose}{What this contributes}
Run independent text or writer passes on multiple processes while restoring page order.


\Q{Method}{Current implementation}
Run independent text or writer passes on multiple processes while restoring page order. The excerpt below is from the current definition; the complete function is in the source snapshot. Its constants and branches, rather than an older signature, govern this path.


\begin{lstlisting}[style=plain]
def _Workers(n):
    if LINE_WORKERS > 0:
        return min(LINE_WORKERS, n)
    return max(1, min(os.cpu_count() or 1, n))
\end{lstlisting}
\Q{Justification}{Interpret the method rather than infer a benchmark}
Explicit fork workers inherit configured models and cap numerical threads to reduce oversubscription. The serial path remains available; startup and IPC still cost time. \cite{pythonmp,numpymatmul}.


\section{\_CapThreads}\label{fn:web-lineworkers:-CapThreads}
\begin{lstlisting}[style=sig]
def _CapThreads()
\end{lstlisting}
\Q{Evidence}{Source and method link}
\path{Software/DemoCode/web/lineworkers.py} line 58. Implementation evidence: \cite{code2026}. Related method, alternatives and library discussion: Section~\ref{method:workers}.


\Q{Purpose}{What this contributes}
Sets thread-count environment variables and tries seven OpenBLAS entry-point spellings (including SciPy-prefixed wheel symbols) found in NumPy shared libraries. It is a best-effort cap, not a portable guarantee for every BLAS/backend. Environment settings after initialisation alone may not reconfigure an already-loaded backend.


\Q{Method}{Current implementation}
Run independent text or writer passes on multiple processes while restoring page order. The excerpt below is from the current definition; the complete function is in the source snapshot. Its constants and branches, rather than an older signature, govern this path.


\begin{lstlisting}[style=plain]
def _CapThreads():
    for v in ("OPENBLAS_NUM_THREADS", "OMP_NUM_THREADS", "MKL_NUM_THREADS"):
        os.environ[v] = "1"
    try:
        import ctypes
        import glob
        import numpy
        root = os.path.dirname(os.path.dirname(os.path.abspath(numpy.__file__)))
        pats = [os.path.join(root, "numpy.libs", "*openblas*.so*"),
                os.path.join(root, "numpy", ".libs", "*openblas*.so*")]
        names = ("scipy_openblas_set_num_threads64_",
                 "scipy_openblas_set_num_threads_64_",
                 "scipy_openblas_set_num_threads",
                 "openblas_set_num_threads64_",
                 "openblas_set_num_threads",
                 "scipy_goto_set_num_threads64_",
                 "goto_set_num_threads")
        for pat in pats:
            for so in glob.glob(pat):
                lib = ctypes.CDLL(so)
                for fn in names:
                    if hasattr(lib, fn):
                        getattr(lib, fn)(1)
                        return True
    except Exception:
        pass
    return False
\end{lstlisting}
\Q{Justification}{Interpret the method rather than infer a benchmark}
Explicit fork workers inherit configured models and cap numerical threads to reduce oversubscription. The serial path remains available; startup and IPC still cost time. \cite{pythonmp,numpymatmul}.


\section{\_Init}\label{fn:web-lineworkers:-Init}
\begin{lstlisting}[style=sig]
def _Init()
\end{lstlisting}
\Q{Evidence}{Source and method link}
\path{Software/DemoCode/web/lineworkers.py} line 91. Implementation evidence: \cite{code2026}. Related method, alternatives and library discussion: Section~\ref{method:workers}.


\Q{Purpose}{What this contributes}
Run independent text or writer passes on multiple processes while restoring page order.


\Q{Method}{Current implementation}
Run independent text or writer passes on multiple processes while restoring page order. The excerpt below is from the current definition; the complete function is in the source snapshot. Its constants and branches, rather than an older signature, govern this path.


\begin{lstlisting}[style=plain]
def _Init():
    _CapThreads()
\end{lstlisting}
\Q{Justification}{Interpret the method rather than infer a benchmark}
Explicit fork workers inherit configured models and cap numerical threads to reduce oversubscription. The serial path remains available; startup and IPC still cost time. \cite{pythonmp,numpymatmul}.


\section{ReadOne}\label{fn:web-lineworkers:ReadOne}
\begin{lstlisting}[style=sig]
def ReadOne(job)
\end{lstlisting}
\Q{Evidence}{Source and method link}
\path{Software/DemoCode/web/lineworkers.py} line 96. Implementation evidence: \cite{code2026}. Related method, alternatives and library discussion: Section~\ref{method:workers}.


\Q{Purpose}{What this contributes}
Run independent text or writer passes on multiple processes while restoring page order.


\Q{Method}{Current implementation}
Run independent text or writer passes on multiple processes while restoring page order. The excerpt below is from the current definition; the complete function is in the source snapshot. Its constants and branches, rather than an older signature, govern this path.


\begin{lstlisting}[style=plain]
def ReadOne(job):
    import numpy as np
    from PIL import Image
    import recognise as REC
    i, crop, flip = job
    img = np.rot90(crop, 2) if flip else crop
    text, score = REC.ReadTextScored(Image.fromarray(img).convert("L"),
                                     _CFG["reader"])
    return i, text, score
\end{lstlisting}
\Q{Justification}{Interpret the method rather than infer a benchmark}
Explicit fork workers inherit configured models and cap numerical threads to reduce oversubscription. The serial path remains available; startup and IPC still cost time. \cite{pythonmp,numpymatmul}.


\section{NameOne}\label{fn:web-lineworkers:NameOne}
\begin{lstlisting}[style=sig]
def NameOne(job)
\end{lstlisting}
\Q{Evidence}{Source and method link}
\path{Software/DemoCode/web/lineworkers.py} line 108. Implementation evidence: \cite{code2026}. Related method, alternatives and library discussion: Section~\ref{method:workers}.


\Q{Purpose}{What this contributes}
Run independent text or writer passes on multiple processes while restoring page order.


\Q{Method}{Current implementation}
Run independent text or writer passes on multiple processes while restoring page order. The excerpt below is from the current definition; the complete function is in the source snapshot. Its constants and branches, rather than an older signature, govern this path.


\begin{lstlisting}[style=plain]
def NameOne(job):
    import numpy as np
    from PIL import Image
    import recognise as REC
    i, crop = job
    _idx, p = REC.ClassifyImage(Image.fromarray(crop).convert("L"),
                                _CFG["namer"], strokeNormalise=_CFG["norm"])
    return i, np.asarray(p, dtype=np.float64)
\end{lstlisting}
\Q{Justification}{Interpret the method rather than infer a benchmark}
Explicit fork workers inherit configured models and cap numerical threads to reduce oversubscription. The serial path remains available; startup and IPC still cost time. \cite{pythonmp,numpymatmul}.


\section{Map}\label{fn:web-lineworkers:Map}
\begin{lstlisting}[style=sig]
def Map(fn, jobs)
\end{lstlisting}
\Q{Evidence}{Source and method link}
\path{Software/DemoCode/web/lineworkers.py} line 124. Implementation evidence: \cite{code2026}. Related method, alternatives and library discussion: Section~\ref{method:workers}.


\Q{Purpose}{What this contributes}
Uses an explicit fork multiprocessing context, creates a pool for the current map and returns results in supplied job order. Failures fall back to serial calls of the same function. Successful model inheritance does not remove crop/result IPC or pool-startup cost, and failure during a serial retry can still propagate.


\Q{Method}{Current implementation}
Run independent text or writer passes on multiple processes while restoring page order. The excerpt below is from the current definition; the complete function is in the source snapshot. Its constants and branches, rather than an older signature, govern this path.


\begin{lstlisting}[style=plain]
def Map(fn, jobs):
    jobs = list(jobs)
    if not jobs:
        return []
    if not CanPool(len(jobs)):
        return [fn(j) for j in jobs]
    try:
        import multiprocessing as mp
        ctx = mp.get_context("fork")
        with ctx.Pool(processes=_Workers(len(jobs)), initializer=_Init) as pool:
            out = pool.map(fn, jobs, chunksize=1)
        return out
    except Exception as e:
        print("[lineworkers] pool unavailable (%s), running in one process"
              % type(e).__name__, file=sys.stderr, flush=True)
        return [fn(j) for j in jobs]
\end{lstlisting}
\Q{Justification}{Interpret the method rather than infer a benchmark}
Explicit fork workers inherit configured models and cap numerical threads to reduce oversubscription. The serial path remains available; startup and IPC still cost time. \cite{pythonmp,numpymatmul}.

